% 编译：xelatex table-algorithm-demo.tex（建议编译两次以更新交叉引用）
% 所有数值均为虚构教学数据；本文件不是任何论文的真实研究结果。
\documentclass[UTF8,a4paper,11pt,fontset=none]{ctexart}
\usepackage[margin=25mm,headheight=15pt]{geometry}
\usepackage{fontspec}
\setmainfont{Liberation Serif}
\setCJKmainfont[BoldFont={Noto Serif CJK SC Bold}]{Noto Serif CJK SC}
\setCJKsansfont[BoldFont={Noto Sans CJK SC Bold}]{Noto Sans CJK SC}
\usepackage{amsmath,amssymb}
\usepackage{booktabs,tabularx,array,threeparttable}
\usepackage{siunitx}
\sisetup{detect-weight=true,mode=text,table-number-alignment=center}
\usepackage[ruled,linesnumbered,noline]{algorithm2e}
\SetAlgorithmName{算法}{算法}{算法目录}
\SetKwInput{KwIn}{输入}
\SetKwInput{KwOut}{输出}
\SetKw{KwReturn}{返回}
\SetAlgoNoLine
\SetAlFnt{\small}
\SetAlgoNlRelativeSize{-1}
\usepackage{caption}
\captionsetup[table]{font=small,labelfont=bf,labelsep=quad,skip=7pt}
\usepackage{fancyhdr}
\pagestyle{fancy}
\fancyhf{}
\fancyhead[L]{\small 数学建模论文\quad 三线表专项优化}
\fancyhead[R]{\small 排版示例}
\fancyfoot[C]{\small\thepage\ /\ 3}
\renewcommand{\headrulewidth}{0.4pt}
\setlength{\parindent}{2em}
\setlength{\parskip}{0.25em}
\linespread{1.22}
\setlength{\tabcolsep}{9pt}
\renewcommand{\arraystretch}{1.18}
\newcolumntype{Y}{>{\raggedright\arraybackslash}X}
\newcommand{\pagetitle}[2]{\noindent{\small #1}\par\vspace{3pt}\noindent{\LARGE\bfseries #2}\par\vspace{7pt}}
\newcommand{\minor}[1]{\par\vspace{6pt}\noindent{\bfseries #1}\par}
\begin{document}

\pagetitle{01\quad 结果对比表}{只突出有依据的最佳值}

同一评价设置下，不同方法的结果如表\ref{tab:results}所示。此处为虚构教学数据，只用于展示数值对齐、局部加粗和表注分工；不包含重复实验统计。

\begin{table}[!htbp]
\centering
\begin{threeparttable}
\caption{不同方法的误差与计算耗时（教学示例）}
\label{tab:results}
\begin{tabular}{@{}l S[table-format=2.2] S[table-format=1.2] S[table-format=1.2]@{}}
\toprule
\textbf{方法} & {\textbf{RMSE}$\downarrow$} & {\textbf{MAE}$\downarrow$} & {\textbf{耗时/s}$\downarrow$}\\
\midrule
基线A    & 12.50          & 9.60          & \bfseries 0.80\\
对照B    & 10.40          & 7.90          & 1.20\\
本文方法 & \bfseries 9.10 & \bfseries 7.10 & 0.95\\
\bottomrule
\end{tabular}
\begin{tablenotes}[flushleft]
\footnotesize
\item[] 注：误差单位与预测量一致；$\downarrow$表示越小越优。粗体表示本例同列最佳报告值，不表示统计显著性。所有数值均为虚构教学数据。
\end{tablenotes}
\end{threeparttable}
\end{table}

本文方法的RMSE为9.10，较对照B降低12.5\%，计算耗时减少20.8\%；基线A的单次耗时最短，为0.80~s。

\minor{为什么没有把“本文方法”整行加粗？}
加粗表达的是同列最佳值，不是方法归属。本文方法的两项误差最小，但最短耗时来自基线A，因此0.95不应加粗。方法名、表头的结构强调与数据标亮应当分开。

\minor{正文解释到这里是否足够？}
这段正文已指出关键误差、相对改善以及时间维度的区别，不必再按行抄录所有数值。是否补充应用意义，取决于任务是否提供明确的误差容限或时间预算；不能自行补造约束。

\minor{什么时候需要更多信息？}
如果真实研究报告多次随机运行，应说明统计量、重复次数和变异来源。如果比较来自不同数据集、设备或预算，应分组或说明条件；不能用一个粗体掩盖不可比性。

\vfill
\noindent{\footnotesize 编排要点：表前短引出；表内精确比较；表注说明读法；表后提炼发现。数值标记只代表声明的比较规则。}

\clearpage
\pagetitle{02\quad 说明类表}{减少散文，不删除必要条件}

表格可以有文字，但文字应当成为便于比较或查阅的字段。下面的处理仅演示表达压缩，不主张任何实际论文都应采用这些操作。

\begin{table}[!htbp]
\centering
\caption{压缩前：操作说明重复使用叙述句}
\label{tab:before}
\small
\begin{tabularx}{\linewidth}{@{}p{25mm}Y@{}}
\toprule
\textbf{处理对象} & \textbf{处理说明}\\
\midrule
序列内部短缺口 & 对两端均有有效观测的序列内部短缺口，采用线性插值的方法对缺失位置进行补全处理。\\
可疑观测 & 根据前文已定义的异常判定规则，对可疑观测进行标记，并进入后续复核步骤。\\
待比较指标 & 根据既定的参数估计范围完成标准化，具体参数来源在实验设置中统一说明。\\
\bottomrule
\end{tabularx}
\end{table}

\begin{table}[!htbp]
\centering
\begin{threeparttable}
\caption{压缩后：按对象、操作和条件组织字段}
\label{tab:after}
\small
\begin{tabularx}{\linewidth}{@{}p{30mm}p{26mm}Y@{}}
\toprule
\textbf{处理对象} & \textbf{操作} & \textbf{关键条件}\\
\midrule
序列内部短缺口 & 线性插值 & 两端观测有效\\
可疑观测 & 标记与复核 & 按前文定义的判定规则\\
待比较指标 & 标准化 & 参数估计范围见实验设置\\
\bottomrule
\end{tabularx}
\begin{tablenotes}[flushleft]
\footnotesize
\item[] 注：这里的“短缺口”、判定规则和参数范围必须在实际论文中有明确的定义和可定位引用；示例不补造这些技术设置。
\end{tablenotes}
\end{threeparttable}
\end{table}

表\ref{tab:after}保留了操作和适用条件，去掉了“采用……的方法”“进行……处理”等重复句法。如果这些内容不承担查阅或对照职责，也可以直接改成正文，而不必保留一张表。

\minor{压缩时保持三个区别}
\textbf{删冗句，不删边界。} 适用条件、来源和关键定义不能随句子一起消失。

\textbf{合理换行，不硬限字数。} 定性说明、规则和符号解释不必全部压成单行。

\textbf{转移有目的，不把冗文搬进表注。} 表注保留阅读所必需的信息；详细原理仍归入正文。

\clearpage
\pagetitle{03\quad 算法三线框}{用控制流展示求解过程}

候选方案的筛选过程见算法\ref{alg:selection}。本例面向给定的有限候选集，采用可行性过滤和目标比较；不是要求所有求解算法都采用这一结构。

\begin{algorithm}[H]
\caption{有限候选集合的可行方案筛选（教学示例）}
\label{alg:selection}
\KwIn{有限候选集$\mathcal{C}$；可行性判定$\operatorname{Feasible}$；可行候选上取有限实值的目标函数$f$（越小越优）。}
\KwOut{选中的可行候选$x_{\mathrm{best}}$，或“无可行候选”状态。}
$x_{\mathrm{best}}\gets\varnothing$；$b\gets+\infty$；$\mathrm{found}\gets\mathrm{false}$\;
\ForEach{$x\in\mathcal{C}$}{
  \If{$\operatorname{Feasible}(x)$}{
    $v\gets f(x)$\;
    \If{$v<b$}{
      $x_{\mathrm{best}}\gets x$；$b\gets v$；$\mathrm{found}\gets\mathrm{true}$\;
    }
  }
}
\If{$\neg\mathrm{found}$}{
  \KwReturn “无可行候选”\;
}
\KwReturn $x_{\mathrm{best}}$\;
\end{algorithm}

算法先过滤不可行候选，再按目标值更新当前方案；遍历结束后，返回给定候选集合内的最佳可行候选。目标值相同时保留先遇到的候选。这里没有声称已找到原始问题的全局最优解。

\minor{算法框与流程图的约定不同}
流程图节点不放公式和具体答案；算法框可以使用必要的变量、赋值、条件和短表达式，以便准确交代计算。长推导放在正文，实际得到的方案和指标放在结果表。

\minor{终止与返回必须明确}
本例的候选集有限，完成遍历即终止，因此不需要虚构“收敛阈值”。无可行候选的状态也有返回路径。迭代算法则应依照真实实现说明更新、停止条件及预算限制。

\minor{算法文字不需要全部加粗}
输入输出标签、控制流关键字提供视觉层次，普通操作保留正常字重；行号用于定位，缩进用于表达嵌套。不要将每一步扩成一段“目的说明”。

\vfill
\noindent{\footnotesize 实现说明：数据表使用booktabs与siunitx；算法使用algorithm2e的ruled风格。实际论文应复用原模板的编号、字体及宏包设置。}
\end{document}
